\begin{figure}[t]
\centering
% Vector drawing generated by tools/figures/relbench_cumulative.py from exp/2026-10-20-relbench/paper-sessions.json; rerun it after the numbers change.
\ifbilingual
\includegraphics[width=\linewidth]{figures/relbench-cumulative-zh}
\else
\includegraphics[width=\linewidth]{figures/relbench-cumulative}
\fi
\caption{\bt{Cumulative tokens, including memory operations, for the six LLM-based methods. Panels share the session axis; shading marks cutoff 1 (learning), followed by cutoffs 2--5 (reuse). Thin rules separate cutoffs.}{六种使用 LLM 的方法的累计 token，含记忆操作。两个面板共用会话轴；阴影标出首截点（学习），其后为截点 2--5（复用）。细线分隔各截点。}}
\label{fig:relbench-cumulative}
\Description{Two vertically stacked line charts, rel-stack and rel-trial, show cumulative tokens in millions over 15 sessions, with one shared legend. The first three sessions are shaded as the initial learning cutoff; sessions four through fifteen are subsequent cutoffs in which memory can be reused.}
\end{figure}
